\documentclass[11pt]{article}

% Keep the PDF trailer identifier stable across isolated reproducible builds.
\ifdefined\pdftrailerid
  \pdftrailerid{0123456789ABCDEF0123456789ABCDEF}
\fi

% Student-facing tutorial: Times-like text/math and an offline, source-only build.
\usepackage[margin=1in]{geometry}
\usepackage{newtxtext,newtxmath}
\usepackage[T1]{fontenc}
\usepackage[utf8]{inputenc}
\usepackage{textcomp}
\usepackage{microtype}
\usepackage{amsmath,mathtools}
\usepackage{booktabs,tabularx,longtable,array}
\usepackage[most]{tcolorbox}
\usepackage{listings}
\usepackage{enumitem}
\usepackage{graphicx}
\usepackage{tikz}
\usepackage{xcolor}
\usepackage{fancyhdr}
\usepackage{lastpage}
\usepackage{hyperref}
\usepackage{xurl}
\usepackage{cleveref}

\definecolor{ink}{HTML}{24292F}
\definecolor{blue}{HTML}{0969DA}
\definecolor{green}{HTML}{1A7F37}
\definecolor{orange}{HTML}{9A6700}
\definecolor{purple}{HTML}{8250DF}
\definecolor{gray}{HTML}{57606A}
\definecolor{codebg}{HTML}{F6F8FA}
\definecolor{codeborder}{HTML}{D0D7DE}
\definecolor{promptbg}{HTML}{EEF6FF}
\definecolor{aibg}{HTML}{F6F0FF}
\definecolor{verifybg}{HTML}{F0FFF4}
\definecolor{checkpointbg}{HTML}{FFF8C5}
\definecolor{warningbg}{HTML}{FFF1E5}
\definecolor{outputbg}{HTML}{F6F8FA}

\hypersetup{colorlinks=true,linkcolor=blue,citecolor=blue,urlcolor=blue,
  pdftitle={AI-Assisted Student Tutorial},pdfauthor={Power Systems Course}}

\lstdefinestyle{tutorialcode}{basicstyle=\ttfamily\small,backgroundcolor=\color{codebg},
  frame=single,rulecolor=\color{codeborder},framerule=0.5pt,columns=fullflexible,
  keepspaces=true,breaklines=true,breakatwhitespace=false,breakautoindent=true,
  postbreak=\mbox{\textcolor{gray}{$\hookrightarrow$\space}},showstringspaces=false,
  xleftmargin=0.5em,xrightmargin=0.5em}
\lstset{style=tutorialcode}
\lstdefinelanguage{CSV}{morestring=[b]",morecomment=[l]{\#}}
\lstdefinelanguage{JSON}{morestring=[b]",morecomment=[l]{//}}

\newtcolorbox{promptbox}{breakable,colback=promptbg,colframe=blue!55!black,
  title=Prompt used in the reviewed example,fonttitle=\bfseries}
\newtcolorbox{aibox}{breakable,colback=aibg,colframe=purple!55!black,
  title=Reviewed AI response (frozen; no live model call),fonttitle=\bfseries}
\newtcolorbox{verifybox}{breakable,colback=verifybg,colframe=green!55!black,
  title=Verify before relying on this, fonttitle=\bfseries}
\newtcolorbox{checkpointbox}{breakable,colback=checkpointbg,colframe=orange!70!black,
  title=Student checkpoint,fonttitle=\bfseries}
\newtcolorbox{warningbox}{breakable,colback=warningbg,colframe=orange!80!black,
  title=Important boundary,fonttitle=\bfseries}
\newtcolorbox{outputbox}{breakable,colback=outputbg,colframe=gray!70!black,
  title=Expected evidence,fonttitle=\bfseries}

\pagestyle{fancy}
\setlength{\headheight}{14pt}
\fancyhf{}
\lhead{AI-Assisted Student Tutorial}
\rhead{ACTIVSg2000 connection planning}
\cfoot{\thepage\ of \pageref{LastPage}}
\newcommand{\code}[1]{\path{#1}}
\newcommand{\MW}{\ensuremath{\,\mathrm{MW}}}
\newcommand{\MWh}{\ensuremath{\,\mathrm{MWh}}}
\emergencystretch=3em

\title{\textbf{AI-Assisted Student Tutorial}\\[0.4em]
  \large ACTIVSg2000 Connection Planning}
\author{Power Systems Course}
\date{October 11, 2026}

\begin{document}
\maketitle

\begin{warningbox}
This tutorial is explanatory. The current student assignment handout and
instructor-released scenario files remain authoritative for rules, sites,
deadlines, scoring, and deliverables. AI text below is representative,
reviewed, and frozen; ordinary PDF and notebook builds do not call a model.
\end{warningbox}

\tableofcontents
\newpage

\section{How to use this tutorial}
Read the short concept sections, run the seven notebooks in order when useful,
and treat each prompt/response pair as a worked example rather than as a new
assignment rule. Start with a small fixture, inspect the evidence, and only
then scale to the released ACTIVSg2000 artifacts. Keep a record of the command,
source version, and output you actually used.

The companion notebooks are \code{01\_dataset\_inspection.ipynb},
\code{02\_network\_visualization.ipynb}, \code{03\_three\_bus\_dcopf.ipynb},
\code{04\_storage\_soc.ipynb}, \code{05\_csv\_validation.ipynb},
\code{06\_student\_preview.ipynb}, and \code{07\_strategy\_ablation.ipynb}.

\section{AI use, privacy, and fact-checking}
The AI-use policy below applies to every prompt, notebook, and submission
workflow in this tutorial.
AI can help explain equations, draft small experiments, and interpret a
traceback. It cannot establish that a bus ID, scenario, score, or evaluator
behavior is current. Ask for verified facts, interpretations, and unknowns in
separate sections. Check generated code on a small fixture, inspect units and
shapes, and compare claims with the handout and repository sources.

Never upload private team plans, credentials, rankings, instructor-only
diagnostics, or unpublished evaluation results to a shared AI tool. Use public
data, synthetic examples, or instructor-approved fixtures. The student preview
is conservative and non-authoritative; it is not an official score.

\begin{checkpointbox}
Before accepting an AI answer, identify its source, reproduce its smallest
claim, check the exact current CSV schema, and write down what remains unknown.
\end{checkpointbox}

\section{Assignment orientation}
The public model is the synthetic Texas ACTIVSg2000 2,000-bus system. The
student-facing task is to prepare a current site-level plan using the released
inputs. The evaluator interprets the plan across one-hour half-year scenarios;
those hours are not rows in the student CSV. Student-owned generation and
storage serve only the submitting team, while the evaluator controls dispatch
and the official result. These statements describe the repository contract;
the authoritative handout controls details not repeated here.

\section{Data and provenance}
Keep the raw public network, the released scenario package, and your derived
figures distinct. The scenario package is an instructor-prepared release built
around the public synthetic model, not a claim about historical ERCOT telemetry.
Record paths and manifest status. A structural network plot is not geographic
evidence unless the input actually supplies coordinates.

\begin{outputbox}
Useful evidence includes file names and shapes, required columns, duplicate
keys, missing-value checks, bus-reference checks, manifest status, and the
exact command used to produce the result.
\end{outputbox}

\section{Dataset inspection}
Begin with \code{01\_dataset\_inspection.ipynb}. Confirm the input files and
their schemas before making a planning claim. A manifest that does not say
\code{validation\_status: passed} must not be described as a validated release.
Use \code{docs/ACTIVSG2000\_SCENARIO\_PACKAGE.md} and the repository README for
provenance and scenario boundaries.

\section{Visualization}
The deterministic visualization notebook produces structural bus/branch,
generator-overlay, and demand-summary views. It should state when geographic
coordinates are unavailable. MATLAB is optional for local exploration; it is
not a substitute for the repository's executable Python checks.

\section{DCOPF}
A DCOPF couples generation, demand, voltage-angle differences, and branch
limits. A useful small example enforces nodal balance and a reference angle
before scaling. For a branch from $i$ to $j$, a common convention is
$f_{ij}=b_{ij}(\theta_i-\theta_j)$, with the sign checked against the chosen
incidence convention. The objective and bounds must be read from the current
inputs; do not invent costs or solver behavior.

\begin{lstlisting}[language=Python,caption={A small verification pattern}]
assert status == "optimal"
assert all(abs(balance) < tolerance for balance in nodal_balances)
assert all(abs(flow) <= limit + tolerance for flow, limit in zip(flows, limits))
\end{lstlisting}

\section{Storage}
The submission expresses storage as power capacity in \MW and energy capacity
in \MWh. The evaluator dispatches storage across the scenario. In a teaching
example, check state transitions, nonnegative state of charge, charge/discharge
power limits, stored-energy bounds, and the terminal state convention. Do not
turn a small illustrative trajectory into a claim about the official score.

\section{CSV submission}
The current canonical site-level CSV has exactly these seven columns:

\begin{center}
\begin{tabularx}{0.98\textwidth}{@{}lX@{}}
\toprule
\textbf{Column} & \textbf{Meaning / unit boundary}\\
\midrule
\code{team\_id}, \code{site\_id}, \code{bus\_id} & identifiers from the released contract\\
\code{connection\_capacity\_mw} & requested connection capacity (MW)\\
\code{generator\_capacity\_mw} & private generation capacity (MW)\\
\code{storage\_power\_mw} & storage power capacity (MW)\\
\code{storage\_energy\_mwh} & storage energy capacity (MWh)\\
\bottomrule
\end{tabularx}
\end{center}

\begin{lstlisting}[language=CSV,caption={Canonical header; verify against the released handout}]
team_id,site_id,bus_id,connection_capacity_mw,generator_capacity_mw,storage_power_mw,storage_energy_mwh
\end{lstlisting}

There are no hourly rows, retired identifier fields, or second submission
schema. Run the repository checker from the repository root:
\begin{lstlisting}[language=bash]
PYTHONPATH=src python scripts/check_submission.py --file team.csv
\end{lstlisting}

\section{Colab and local preview}
The notebooks can be opened locally or in Colab after providing the repository
and its versioned inputs. Run the validator before the preview, save its output,
and label preview metrics as non-authoritative. The preview helps diagnose
format and modeling issues; only the official evaluation determines the score.

\section{Strategy experiments}
Use \code{07\_strategy\_ablation.ipynb} to vary one factor at a time on a
controlled fixture. Record candidate parameters, validation status, preview
metrics, and the figure produced. An ablation can reveal sensitivity; it does
not prove an optimal strategy and must not include another team's private data.

\section{Debugging and reproducibility}
When a cell fails, preserve the traceback, input path, package version, and
smallest reproducer. Separate a path/schema problem from a numerical or solver
problem. Re-run after fixing one layer, and keep the exact command in your
notes. The repository scripts provide the stable entry points for notebook
execution and tutorial validation.

From the repository root, the validation sequence is:
\begin{lstlisting}[language=bash]
python docs/student_tutorial/scripts/validate_prompt_registry.py
python docs/student_tutorial/scripts/validate_transcripts.py
python docs/student_tutorial/scripts/execute_notebooks.py \
  --input-dir docs/student_tutorial/notebooks \
  --output-dir docs/student_tutorial/outputs/notebook_runs
python docs/student_tutorial/scripts/validate_tutorial_examples.py
\end{lstlisting}

\section{Final checklist}
\begin{itemize}[leftmargin=*]
  \item I used the authoritative handout and released scenario/catalog inputs.
  \item I verified the exact seven-column site-level CSV schema.
  \item I did not add hourly rows or unsupported fields to the submission.
  \item I checked units, bus IDs, file paths, and validator output.
  \item I treated the preview as non-authoritative.
  \item I protected private team information and recorded provenance.
  \item I can explain every generated code fragment I submit.
\end{itemize}

\section{Reviewed prompt and transcript examples}
The following sections are generated from the frozen, fact-checked JSON
transcripts. Run \code{python docs/student\_tutorial/scripts/build\_latex\_inputs.py}
before compiling from a clean checkout. The generator escapes ordinary prose
and routes fenced code to \code{listings}.

\IfFileExists{generated/transcripts.tex}{\input{generated/transcripts.tex}}{%
  \begin{warningbox}
  Transcript includes are absent. Run \code{build\_latex\_inputs.py} before
  compiling the complete tutorial.
  \end{warningbox}}

\section{References and appendices}
The repository's assignment handout, student manual, scenario-package guide,
and current checker instructions are the first sources to consult. The
background tools listed below document formats and reproducible notebook
execution, not the assignment contract.
\nocite{activsg2000,matpower,colab,tcolorbox,nbclient,jupyterbook,quarto,nbgrader}
\bibliographystyle{plain}
\bibliography{references}

\end{document}
